\scanpage{105}

façon unique et assure son existence et unicité globales. On \deleted{\unsure{a donc}} trouve alors que sur $\underset{\substack{\Vert\\V\times1}}{E}\deleted{\unsure{=V}}\inserted{\subset\mathscr E}$ la fonction induite par $Q$ est de la forme $q(x)+1$, où $q$ est la forme quadratique \inserted{lisse} induite sur $V$ par $Q$. Donc la donnée de $\Sigma$ dans $E$ (\unsure{quadrique|épithète} non paraboloïde) équivaut à celle des

a) Section $o$ de $E$ (le \uline{centre} de $\Sigma$)

b) forme quadratique \inserted{lisse} sur $V$ (fibré vectoriel des translations)

en associant à ces données

\[
\Sigma=\unsure{\mathcal V|V|\check{\mathbb V}}(q+1)=\{x\in\Gamma(-,V)\mid q(x)=-1\}
\]

(en identifiant $E$ à $V$ grâce à $o$).

Ces développements sont valables pour tt rang de $E$, en prenant soin, quand celui-ci est impair, \inserted{(comme p.ex. rg 3 comme ici)} de se restreindre au cas de \deleted{\unsure{la}} car. rés. $\ne2$.

Considérons \deleted{\unsure{$\delta$}}

\[
\delta'(q^{\unsure{\prime|\ast}})\in(\det\check V)^{\otimes2}\qquad\text{base de }(\det\check V)^{\otimes2},
\]

% archive sheet number: 104
